\chapter{Project Management}
\label{ch:mgmt}

\section{Development Approach}
The project follows an iterative process. Each iteration ends with a working, tested increment and a review of what to keep, change or drop; the floor-distance prototype, for example, was dropped after iteration~1 when moving-camera footage showed it could not work. Work is split into owned work items (for example B1 labels, C3 dispatch, M-UI-P3 dashboard connection), each developed on a feature branch, reviewed and merged into a shared integration branch. Status, contracts and open items are kept in tracked planning documents in the repository so every member works from the same information.

\section{Timeline and Milestones}
\begin{longtable}{P{0.2\textwidth}P{0.72\textwidth}}
\caption{Milestones}\label{tab:milestones}\\
\toprule
Date & Milestone\\
\midrule
\endfirsthead
\toprule
Date & Milestone\\
\midrule
\endhead
\bottomrule
\endfoot
26 Aug 2026 & First commit to the repository\\
11 Sep 2026 & Proposal defence (feedback in Appendix~\ref{app:feedback})\\
12 Oct 2026 & Supervisor accepts the team's 70-video dataset\\
21 Oct 2026 & Mid evaluation: vision prototype and web application (iteration~1)\\
24 Oct 2026 & Model consolidation: COCO detector tracks everyone, role detector only labels\\
3 Nov 2026 & Iteration~2 begins with the team's 16-week plan\\
5--8 Nov 2026 & 25 clips labelled; benchmarks on tuning and held-out sets\\
10 Nov 2026 & Agent foundation: typed tools, scoped registry, orchestration graph\\
11 Nov 2026 & Move to PostgreSQL; dashboard connected to the backend\\
12 Nov 2026 & Business agents, signal policy, occupancy categories, digest, clip review, CI workflow\\
13 Nov 2026 & Security fixes; direct dispatch connected end to end\\
14 Nov 2026 & Role-weight decision; acceptance rehearsal on the hosted database\\
15--18 Nov 2026 & Map registration, directional footfall, frontend security, production audit\\
17 Nov 2026 & Benchmarks rerun with fitted weights; 276 backend tests\\
25--26 Nov 2026 & Aisle relay demonstration; dispatch skips employees in an interaction\\
\end{longtable}

\section{Team Responsibilities}
\begin{table}[H]
\centering
\caption{Ownership of components}\label{tab:ownership}
\begin{tabular}{P{0.2\textwidth}P{0.72\textwidth}}
\toprule
Member & Responsibilities\\
\midrule
Moiz Ansari (23I-0523) & Manager frontend and dashboard binding; agent foundation (typed tools, scoped registry, orchestration graph, operational adapter); Ask route; footfall counter; PostgreSQL migration\\
Saqib Atif (23I-0769) & Vision pipeline, labelling and benchmarks; event store and ingestion; dispatch and voice; map registration; security fixes; aisle relay demonstration\\
Ruhab Ahmad (23I-0559) & Employee frontend; Customer Behavior and Employee Activity agents; signal policy; occupancy categories; digest; clip review; golden questions; acceptance tests; security review; CI workflow; role-weight decision; documentation and this report\\
\bottomrule
\end{tabular}
\end{table}

\section{Version Control Evidence}
The repository \path{muhammad-moiz-ansari/retail-store-analysis} has 339 commits on the branch described here, from 26 August to 26 November 2026. Git records them under five author identities: Moiz Ansari (136), \texttt{msaqibatifj} (128) and Muhammad Saqib Atif (2) for Saqib, Ruhab Ahmad (62), and \texttt{GulsherKhan-04} (11). The last identity holds Saqib's labelling and benchmark commits of 5--8 November, which were recorded under the wrong identity and re-authored to him on the shared integration branch. Commit counts measure activity, not the size or value of each contribution.

\section{Risk Management}
\begin{longtable}{P{0.22\textwidth}P{0.34\textwidth}P{0.36\textwidth}}
\caption{Risks and responses}\label{tab:risks}\\
\toprule
Risk & Effect & Response\\
\midrule
\endfirsthead
\toprule
Risk & Effect & Response\\
\midrule
\endhead
\bottomrule
\endfoot
No store would share footage & No data to train or evaluate & Team recorded its own 70-video dataset, accepted by the supervisor\\
No GPU & Slow processing & Small models, OpenVINO, frame sampling and role labelling only for new tracks; 0.60$\times$ real time achieved\\
Moving cameras & Floor mapping failed & Floor prototype descoped; body-height distance used\\
Small labelled set & Unreliable accuracy figures & Held-out set, bootstrap intervals, error analysis, blind re-check planned\\
Language model invents figures & Wrong answers to managers & Claims must cite tool metrics; grounding and narrative checks; model used only for routing\\
Duplicate or wrong dispatch & Staff sent twice or a zone left empty & Unique event key; eligibility protects minimum staffing; scenario tests\\
Exposed secrets & Forged sessions & Startup refuses default secrets; secrets server-side; keys exposed earlier are being rotated\\
Store map not supplied & Map-cell occupancy cannot be shown & Map registration tested on synthetic geometry; real map requested for FYP-2\\
\end{longtable}

\section{Tools and Generative-AI Use}
Development used Git and GitHub, GitHub Actions, CVAT for labelling, Node.js, a Python virtual environment and PostgreSQL. Generative-AI assistants were also used, as disclosed in Table~\ref{tab:ai-use}. Every AI-assisted change was checked by tests and by the member who owns the affected component, and each member must be able to explain the code and text in their area.

\begin{table}[H]
\centering
\caption{Generative-AI use and how it was checked}\label{tab:ai-use}
\begin{tabular}{P{0.15\textwidth}P{0.4\textwidth}P{0.37\textwidth}}
\toprule
Tool & Used for & Check\\
\midrule
Claude Code & Implementation support for agents, tests, acceptance rehearsal, report drafting and figure redrawing & Automated tests; benchmark reruns; every report number traced to its source file; figures checked against the code\\
Codex & Implementation support for the descoped floor prototype and documentation & Automated tests and output inspection\\
Claude & General development questions & Review by the owning member\\
\bottomrule
\end{tabular}
\end{table}
